%\input ../util/bookmacros
%\input ../util/docparams

\chapter{Introduction to Difference Equations}

\section{Motivational Problem}
If a person buys a house they will almost surely need a loan from a bank 
to pay for it. Once you borrow the money, the bank will have you spread 
the payments out over $15$ to $30$ years. You're not simply paying
back the money to the bank; it charges interest on the use of its
money. So, each payment can be thought of 
as consisting of two pieces: a payment for the principal borrowed and a 
second piece which is interest on the outstanding debt. We now describe what 
happens at any given month. 

At month $n$ let $x_n$ be
the outstanding debt we owe to the bank. The bank has decided that 
at the end of each month we pay $p$ dollars for the loan. Also, 
during the month, we are charged interest (for that month) on the 
outstanding debt as 
of the beginning of that month. The rate of interest is given as an annual interest rate (AIR) 
and so must be divided by $12$ to get the monthly interest rate, 
$r = {\rm AIR}/12$. At the beginning (month zero) the outstanding 
debt is known, it is just the loan amount; we'll call it $L$.

We would like to know what the outstanding debt is at any given month. 
Our strategy is similar to what one would do in an algebra class: rather than 
writing down $x = $ {\it answer\/}, one first found a 
relation that $x$ satisfied and 
then solved for $x$. We will do the same thing here. Rather than write 
down a formula for the outstanding debt at month $n$, $x_n$, we first 
find a relation that $x_n$ satisfies and then solve for $x_n$. Such a relation 
can often be stated as a change or a {\it \index{difference}\/} from one time 
step to the next. This change is often easy to describe. For the problem at hand, the change in debt from one 
month to the next is the amount of new debt accumulated by interest less the 
\index{monthly payment}.
Mathematically, this can be written as:
$$
\eqalign{
\obrace{x_{n+1}}{Debt at month $n{+}1$} & = \obrace{x_n}{Debt at month $n$}
+ \obrace{r x_n}{Interest on debt over the month} -
\obrace{p}{Monthly payment} \cr
x_0 & = L \cr
}
$$
This can be rewritten as:
$$
\eqalign{
{x_{n+1}} & = (1 + r) x_n - p \cr
x_0 & = L \cr
}
$$
\section{Comparison with the Solution Strategy of Algebra}
Recall that in high school algebra one solved ``word'' problems in a two step procedure:
\beginEnum
	\enum{Label the unknown {\it number\/}, $x$, and involve it in an algebraic equation.}
	\enum{Treat the equation as a ``balance scale'' and apply the operations of 
			addition/subtraction and/or multiplication/division equally on 
			both sides in an attempt to isolate $x$ on one side 
			-- its value is then on the other side of the equation.}
\endEnum
Here, the unknown labeled, $x$, is not a number, it is 
a {\it \index{sequence}\/} of numbers; so, what does it mean to solve for $x$?
Ideally, we would like a {\it formula\/} for the $n^{\rm th}$ element of 
the sequence $x$ in terms of this $n$. Then, when 
someone asked for the debt at month $17$, we could just plug $17$ 
into this formula and out would pop the debt for that month.

\section{Basic Solution Strategy}
The equation we wrote down for the debt as a function of the month 
was an example of a {\it \index{difference}} equation: each term in the sequence is a function of previous 
terms. The difference and summation calculus of the 
first two chapters
will play an important role in solving such equations. 
In the equation above, a given element of the sequence is a function of the previous element.
Such an equation is called a first order {\it \index{difference equation}}.
The form of this relation between previous elements 
determines how easy it is to solve, depending 
on what the meaning of 
``solve'' is. Before we start, let us be clear: The solution to any difference 
equation is not a number; that is, the unknown we are looking to 
find is an entire {\it \index{sequence}\/} -- which we have said is a {\it function\/} on the integers.%
\footnote{\kern 0.5pt \raise 0.5ex \hbox{\dag}}{Described in the sub-section ``Sequences as Functions'' in the first chapter.}
Again, we are doing something very different from high school algebra -- we are solving for a {\it function\/}.

Let's start with a simpler difference equation. In fact, let's start with 
the simplest first order difference equation relating a given value of 
a sequence with its previous value:
$$
\eqalign{
x_{n+1} & = x_n \cr
x_0 & = L \cr
}
$$
We can solve this difference equation by using the equation to write 
down the solution term by term. That is, we know $x_0 = L$. 
The first equation tells us that $x_1 = x_0$; and since we know $x_0$, we 
know $x_1$. Continuing, we have $x_2 = x_1$; and since we know $x_1$, 
we know $x_2$, it is just $x_1 = x_0 = L$. We can continue and produce 
each value of the sequence. It is tedious, but effective. Every \index{difference 
equation} can be solved by doing this sort of computation. In theory, any 
element of the sequence could be produced this way. In practice, one 
would compute as many elements of the sequence as one needed.

A better solution though, would be to write down a {\it formula\/} for 
any given element of the sequence. In our present example, we find 
that the formula 
for the $n^{\rm th}$ element of the \index{sequence} is $x_n = L$. Having a formula 
means one can, on request, produce any element of the sequence without 
having to compute {\it any\/} of the previous elements. 
The question is 
\goodbreak
``When presented with a difference equation, can we write down a formula 
for the $n^{\rm th}$ term of the sequence satisfying such an equation?''
The answer is no in general. We will, however, be able to find a solution (in this sense) to our motivational problem.

The problem with 
finding a formula, in general, is that the $n^{\rm th}$ value of the sequence 
which satisfies such an equation is tangled up with the previous values. 
Because of this, it is not a simple matter to write down a formula 
for the $n^{\rm th}$ value. We introduce a technique 
which will allow us to solve a larger array of equations.
Consider the previous simple \index{difference equation}; 
we know by the \index{Fundamental Theorem of Difference Calculus} that 
it is easy to compute the sum of \index{telescoping difference}s. Therefore, 
we may rewrite the last equation as:
$$
\eqalign{
x_{n+1} - x_n & = 0 \cr
x_0 & = L \cr
}
$$
This is
$$
\eqalign{
\Delta x_n & = 0 \cr
x_0 & = L \cr
}
$$
Summing both sides of the first equation from $0$ to $N-1$ we have:
$$
\sum_{n=0}^{N-1} \Delta x_n = \sum_{n=0}^{N-1} 0 
$$
After summing (using the FTDC on the left hand side):
$$
x_N - x_0 = 0
$$
Or, $x_N = L$.
We were able to untangle the connections of the sequence by summing them out.
This left an element of our choosing, $N$, and the $0^{\rm th}$ term plus
whatever the sum was on the right hand side of the equation.
Notice that this is the same game we have been playing since algebra: we
isolated the unknown under the operator $\Delta$, then inverted $\Delta$ by
applying its inverse, summation -- the {\it isolate and invert\/} move of the preface.

If we modify the difference equation so that it becomes:
$$
\eqalign{
x_{n+1} & = x_n + b_n \cr
x_0 & = L \cr
}
$$
This is
$$
\eqalign{
\Delta x_n & = b_n \cr
x_0 & = L \cr
}
$$
Summing both sides of the first equation from $0$ to $N-1$ gives:
$$
\sum_{n=0}^{N-1} \Delta x_n = \sum_{n=0}^{N-1} b_n 
$$
After summing (using the FTDC):
$$
x_N - x_0 = \sum_{n=0}^{N-1} b_n
$$
Or, $x_N = L + \sum_{n=0}^{N-1} b_n$. If $b_n = b$, that is, 
$b_n$ is constant, then $x_N = L + b N$.

\subsection{Solution of the Motivational Problem}
We are moving toward the solution of our motivational problem. To get a step 
closer, consider the following difference equation:
$$
\eqalign{
x_{n+1} & = a\, x_n \cr
x_0 & = L \cr
}
$$
A solution to this \index{difference equation} is a sequence which has $L$ for 
its initial value and then multiplies the current value by $a$ to get the next.
It is not too hard to guess that the solution%
\footnote{\kern 0.5pt \raise 0.5ex \hbox{\dag}}%
{The pattern emerges quickly. $x_1 = a \,x_0 = a \,L$; $x_2 = a \,x_1 = a \, a \, L = a^2 \, L$; $x_3 = a \, x_2 = a \, a^2 L =
a^3\, L$.} is $x_n = a^n L$.
Let us check this is indeed a solution%
\footnote{\kern 0.5pt \raise 0.5ex \hbox{\ddag}}{In fact, it is {\it the\/} solution.}. 

We need to check two things. First, that 
$x_0 = L$. This is so, since $x_0 = a^0 L = 1 * L = L$. We also need 
to check $x_{n+1} = a\, x_n$. But $x_{n+1} = a^{n+1}L = a (a^n L) 
= a\, x_n$. So this is a solution. Remember, a solution to a difference 
equation is a {\it sequence\/}, not a number.

Now let's look at an equation that has the same form as our original problem:
$$
\eqalign{
x_{n+1} & = a\, x_n + b_n \cr
x_0 & = L \cr
}
$$
It is not as easy to guess a solution to this.
Can we also use the differencing technique we used earlier? Well, if we rewrite 
the equation so that the left hand side is a difference as before, then:
$$
\eqalign{
x_{n+1} - x_n & = (a-1)x_n + b_n \cr
x_0 & = L \cr
}
$$
This does not help us untangle the sequence of unknowns, $x_n$. If we try to sum 
this equation, the left hand side reduces to $x_N$ and $x_0$. However, the right hand 
side contains the unknowns for which we are trying 
to solve. For this technique to work we need to get all of the $x_n$'s 
on the left as part of {\it some\/} telescoping difference (not necessarily a difference 
in $x_n$ alone), leaving the right hand side containing 
the rest of the {\it known\/} quantities. Summing the left hand side will reduce 
to two terms; from there we can solve for $x_N$. To start, let's 
rewrite the equation as a difference of the unknowns (not yet a \index{telescoping difference}):
$$
\eqalign{
x_{n+1} - a\, x_n & = b_n \cr
x_0 & = L \cr
}
$$
The key idea is to multiply both sides of the first equation by a known \index{sequence} 
so that the left hand side is a quantity which can be represented as 
a {\it telescoping\/} difference. Given a sequence 
$\{y_n\}_{n=0}^\infty$, multiply the first equation by $y_{n+1}$. This gives:
$$
\eqalign{
y_{n+1} x_{n+1} - (a\, y_{n+1}) x_n & = b_n y_{n+1} \cr
x_0 & = L \cr
}
$$
If we chose the sequence $y$ so that $a\, y_{n+1} = y_n$, then we would have 
a \index{telescoping difference} on the left hand side 
$$
\eqalign{
y_{n+1} x_{n+1} - y_n x_n & = b_n y_{n+1} \cr
x_0 & = L \cr
}
$$
Such a multiplier sequence $y$ is called a {\it \index{summing factor}\/} -- and
when we solve differential equations in the next chapter, we will meet its
continuous twin, the {\it integrating factor\/}.

Summing from $0$ to $N-1$ gives us $y_N x_N - y_0 x_0 = \sum_{n=0}^{N-1} b_n y_{n+1} $. Solving for $x_N$ 
gives
$$
x_N = \biggl( {y_0 L + \sum_{n=0}^{N-1} b_n y_{n+1} \over y_N}\biggr)
$$
So, we're done provided we know 
how to find the sequence $y$. But the sequence $y$ 
satisfies its own \index{difference equation}: $y_{n+1} = a^{-1} y_n$. From what we discovered above, the solution 
is $y_n = a^{-n} y_0$. There are no further restrictions so we may choose 
$y_0$ to be anything; let's set it to $1$. Therefore, we may write down 
a solution for the $N^{\rm th}$ element of the sequence $x$ to be
$$
x_N = a^N L + a^N \sum_{n=0}^{N-1} a^{-(n+1)} b_n
$$
When $b_n$ has the same value, say $b$, for all $n$ this formula simplifies to
$$
x_N = a^N L + b\, a^{N-1} \sum_{n=0}^{N-1} a^{-n}
$$
Which simplifies to 
$$
x_N = a^N L + b\, a^{N-1} \biggl( {a^{-N} - 1 \over a^{-1} - 1} \biggr)
$$ 
provided $a \ne 1$ (we used the following fact about \index{geometric series}: $\sum_{n=0}^{N-1} z^n = {z^{N} - 1 \over z - 1}$, for $z \ne 1$;
used above with $z = a^{-1}$).
This can also be written as:
$$
x_N = a^N L + b\, \biggl( {1 - a^{N} \over 1 - a} \biggr)
$$
The two terms of this solution each tell a story. The first term, $a^N L$,
solves the equation with no payments ($b = 0$) and carries the initial
condition; the second term solves the full equation starting from zero.
Because the equation is {\it linear\/}, the two pieces simply add -- this is
the {\it exploit linearity\/} move ({\it \index{superposition}\/}).

Applying this formula to the motivational problem with $a = (1 + r)$, and
$b = -p$ yields a formula for the debt owed at month $N$:\footnote{\kern 0.5pt \raise 0.5ex \hbox{\ddag}}{Note $a \neq 1$ so we may use the just derived formula.}
$$
x_N = (1+r)^N L - p \biggl( {1 - (1+r)^{N}  \over -r} \biggr) 
$$
which is
$$
x_N = (1+r)^N L - p \biggl( {(1+r)^{N} - 1  \over r} \biggr) 
$$

\example{If a $30$ year loan is taken out on \$$100,000$ at an 
annual interest rate of $6$\% with monthly payment 
\$$600$, what is the outstanding debt after $5$ years?}

The monthly rate is $r = 0.06 / 12 = 0.005$, $L = 100,000$, $p = 600$, and 
$N = 5 * 12 = 60$. Plugging these values into the solution to our difference 
equation gives
$$
x_{60} = (1.005)^{60} * 100000 - 600 * 
\biggl({1.005^{60} - 1 \over 0.005} \biggr)
$$
After evaluating the expression on the right hand side we obtain
$x_{60} = 93,023$. 
This is the outstanding debt after $5$ years of monthly payments.
During this time you will have paid $5 * 12 * 600 = 36,000$
dollars to the bank.

\exercise{Solve the difference equation $x_{n+1} = 2\, x_n + 1$, $x_0 = 0$,
using the formula of this section. (The answer, $x_N = 2^N - 1$, counts the
moves needed to solve the Tower of Hanoi puzzle with $N$ disks.)}

\exercise{You deposit $L$ dollars in a savings account earning monthly interest
rate $r$, and add $p$ dollars at the end of each month. Write down the
difference equation satisfied by the balance and solve it. How does the
solution compare, term by term, with the loan formula above?}


\section{Summary}
We are able to solve a simple \index{difference equation} by first 
rewriting the tangling of $x$'s as a \index{telescoping difference} and then use 
the \index{Fundamental Theorem of Difference Calculus} to ``sum away'' the 
tangles. This leaves one element of the sequence, $x_N$, on one 
side of the equation and the sum of 
the known values on the other side of the equation.

\section{What's Left}
A difference equation is effectively any equation which relates a given element of a sequence with previous values of the 
sequence. Obviously, they can
be much more complicated than the ones we've examined. A careful study of the subject classifies difference 
equations into 
various categories and then provides recipes for solutions in some cases. The main solution technique is to
``sum away'' the differences; that is, use the \index{Fundamental Theorem of Difference Calculus}.
 
\subsection{Fibonacci Sequence}
As an example we consider a more complicated difference equation.
The Fibonacci sequence occurs in many contexts, one of which is the modeling of rabbit populations. The
sequence is defined through a recurrence relation (a difference equation) given by:%
\footnote{\kern 0.5pt \raise 0.5ex \hbox{\dag}}{Some texts index the sequence
differently, starting with $F_0 = 0$ and $F_1 = 1$; this produces the same
list of numbers shifted by one place.}
$$
\eqalign{
F_{n+2} & = F_{n+1} + F_{n}, \quad n \ge 0 \cr
F_0 & = 1 \cr
F_1 & = 1 \cr
}
$$
This is categorized as a second order difference equation because a given element in the sequence is a
function of the previous two values.

How might one solve it? A change of representation does the trick: try a
solution of the form $x_n = \lambda^n$. Substituting into the recurrence and
dividing by $\lambda^n$ forces $\lambda^2 = \lambda + 1$, a quadratic
whose roots, $(1 \pm \sqrt{5})/2$, involve the famous {\it \index{golden ratio}\/}.
Because the recurrence is {\it linear\/}, any combination
$A\,\lambda_1^n + B\,\lambda_2^n$ of the two solutions is again a solution,
and the two initial values pin down $A$ and $B$.

\exercise{Carry this out: find $A$ and $B$ so that
$F_n = A\,\lambda_1^n + B\,\lambda_2^n$ satisfies $F_0 = F_1 = 1$, and thereby
obtain a formula for the $n^{\rm th}$ Fibonacci number. Check your formula by
computing $F_2$ and $F_3$.}

